\section{Long CoT visual: llevar el pensamiento lento al espacio visual}

La sección anterior trató el razonamiento en lenguaje; esta sección vuelve a lo multimodal. Después de estudiar o1, Zhang Xiangyu (张祥雨) volvió a pensar el problema de la poca controlabilidad de la generación visual: muchas tareas de generación visual también superan, en esencia, la complejidad de un razonamiento de un solo paso. Cuando una persona dibuja o comprende relaciones espaciales, no genera una imagen completa a partir de ruido de una sola vez; primero comprende la escena, localiza las partes, revisa las relaciones y después corrige paso a paso. Si lo multimodal solo hace generación de una sola vez, viola con facilidad las restricciones geométricas y físicas.

\begin{figure}[H]
\centering
\includegraphics[width=0.96\textwidth]{visual-long-cot.png}
\caption{Long CoT en el espacio visual: pensar despacio en el espacio visual y trasladar el patrón de reflexión al razonamiento multimodal. Diagrama conceptual propio, elaborado a partir de la conversación de 01:10:57--01:24:07.}
\end{figure}

\begin{knowledgebox}{Lectura de la figura: el CoT visual no consiste en alargar el razonamiento escrito en chino}
La figura va de Image/Video a Visual Actions, Reason, Reflect y Answer. Lo central del Long CoT visual es que el modelo realice acciones en el espacio visual: observar una región, anotar, volver a mirar, comparar, generar borradores auxiliares y corregir relaciones espaciales, en lugar de solo producir explicaciones de texto más largas.
\end{knowledgebox}

\subsection{Por qué la generación también necesita CoT}

Esta sección conecta la comprensión visual con la generación visual. La ruta de Zhang Xiangyu es empezar por el CoT en la comprensión visual, porque en la comprensión es más fácil obtener retroalimentación verificable; pero el objetivo de más largo plazo es que la generación también incorpore una cadena de pensamiento. La generación muy controlable no se dibuja de una sola vez: primero se comprenden las restricciones semánticas y después se aproxima el objetivo mediante generación local, revisión y modificación. En videos complejos esto es todavía más importante, porque la continuidad temporal, la causalidad física y la consistencia de identidad deben mantenerse entre cuadros.

\begin{knowledgebox}{Términos clave: espacio de acciones visuales}
\noindent\begin{tabularx}{\textwidth}{p{0.20\textwidth}p{0.34\textwidth}X}
\toprule
Acción & Significado & Por qué es útil \\
\midrule
Observación local & Mirar solo una parte de la imagen o del video & Reduce la interferencia del contexto y mejora el juicio sobre los detalles. \\
Anotar/trazar recuadros & Generar estructuras intermedias en el espacio visual & Hace explícitas las relaciones espaciales implícitas. \\
Volver a mirar/comparar & Contrastar distintas regiones o momentos & Detecta inconsistencias y errores físicos. \\
Generar borradores & Generar primero resultados intermedios que se puedan revisar & Permite que el proceso de generación reciba retroalimentación y se corrija. \\
Reflexión y corrección & Rehacer el camino tras detectar un error & Convierte la tarea visual en un proceso de búsqueda. \\
\bottomrule
\end{tabularx}
\end{knowledgebox}

\subsection{Dos frentes: corpus de preentrenamiento y espacio de acciones}

Lo anterior mostró que el Long CoT visual necesita acciones; esta sección concentra la ruta en dos frentes. La ruta futura de lo multimodal no es una sola línea. El primer frente es ampliar el corpus de preentrenamiento: más datos de alta calidad de imagen y texto, de video, de procesos de generación, de procesos de interacción y de razonamiento visual. El segundo frente es ampliar el espacio de acciones: que el modelo no solo pueda ver y hablar, sino también observar regiones, dibujar, generar, revisar, llamar herramientas, llamar submodelos y recibir retroalimentación. Solo si los datos y el espacio de acciones se amplían juntos, el pensamiento visual tendrá la oportunidad de acercarse al salto que dio el razonamiento en lenguaje.

\begin{figure}[H]
\centering
\includegraphics[width=0.94\textwidth]{two-leg-roadmap.png}
\caption{Los dos frentes futuros de lo multimodal: ampliar el corpus de preentrenamiento y, al mismo tiempo, ampliar el espacio de acciones. Diagrama conceptual propio, elaborado a partir de la conversación de 01:31:31--01:45:42.}
\end{figure}

\begin{knowledgebox}{Lectura de la figura: los datos resuelven «haberlo visto»; el espacio de acciones resuelve «saber hacerlo»}
A la izquierda, Pretraining Data permite que el modelo vea más procesos visuales y de generación; a la derecha, Action Space le da al modelo más acciones ejecutables dentro de la tarea. Es muy probable que el «momento GPT-4» de lo multimodal requiera que ambas cosas ocurran a la vez, en lugar de seguir solo acumulando pares de imagen y texto.
\end{knowledgebox}

\subsection{Línea de tiempo de la ruta de investigación}

Esta sección condensa los conceptos anteriores en una línea de tiempo, para que el lector entienda este episodio desde la historia de la investigación y no desde las palabras de moda. El relato de Zhang Xiangyu pasa aproximadamente por cuatro etapas: la primera es el CV scaling, hacer más profundos y más grandes modelos como ResNet; la segunda es la CV aprendiendo de NLP, con el intento de trasladar a las imágenes el Transformer, BERT/GPT, MIM y las ideas autorregresivas; la tercera es el tropiezo de la integración multimodal, al descubrir que la generación y la comprensión no se fortalecían de verdad entre sí; la cuarta es volver a entender el razonamiento después de o1 y llevar de regreso a la visión el RL, el CoT, la reflexión y el espacio de acciones.

\noindent\begin{tabularx}{\textwidth}{p{0.16\textwidth}p{0.31\textwidth}X}
\toprule
Etapa & Acción central & Lección obtenida \\
\midrule
2012--2016 & ImageNet, AlexNet y ResNet demuestran el model scaling en visión & Ampliar a la vez datos, cómputo y modelo abre el reconocimiento visual. \\
2018--2022 & ViT, iGPT, MIM, aprendizaje contrastivo y otros aprenden de NLP & La transferencia de arquitecturas funciona, pero la función objetivo y las propiedades de los datos no se resolvieron de fondo. \\
2023--2024 & Entrenar grandes modelos multimodales e intentar integrar generación y comprensión & La comprensión y la generación se fortalecen cada una por su lado, pero no se fusionan automáticamente. \\
Después de o1 & Volver a ver lo multimodal desde el RL, el CoT, la reflexión y el espacio de acciones & La visión también necesita pensamiento lento y procesos intermedios en los que se pueda buscar. \\
Próximos dos años & Long context, aprendizaje en línea, aprendizaje autónomo y modelos del mundo & El momento GPT-4 podría venir de la retroalimentación del entorno y de la colaboración entre sistemas. \\
\bottomrule
\end{tabularx}

\begin{importantbox}{El docente enfatiza: la desorientación misma es evidencia sobre la ruta}
En la entrevista aparecen una y otra vez expresiones como «pasé más de medio año muy desorientado» o «cuando salió o1 lo vi claro». Esa desorientación no es ruido, sino evidencia sobre la ruta técnica: si una dirección sigue sin lograr 1+1>2 aunque se agreguen datos, se agrande el modelo o se conecten módulos, lo que falta es un nuevo objetivo de entrenamiento o un nuevo espacio de acciones.
\end{importantbox}

\subsection{Resumen de la sección}

El Long CoT visual es la idea multimodal nueva más importante de este episodio: trasladar a la visión las ideas de reflexión y de espacio de acciones de o1, para que el modelo piense despacio en el espacio visual. No se trata de alargar la cadena de texto, sino de darle al modelo multimodal más acciones visuales verificables, en las que se pueda buscar y que se puedan corregir.
